% =====================================================================
%  Betriebsanweisung Dübelfräse Festool DOMINO DF 500 RQ-Plus
%  Eigenes PDF: Betriebsanweisung_Duebelfraese.tex, Seite in der Einweisung
%  über \BAseite. Layout: fablab-document/fablab-ba.sty
%
%  Lizenz-Hinweis: Text vollständig selbst formuliert, KEINE Passagen und
%  KEINE Abbildungen aus der Festool-Betriebsanleitung übernehmen! Sicherheits-
%  regeln als solche (Fakten) sind frei, der Wortlaut und die Bilder von
%  Festool sind urheberrechtlich geschützt. Nur ISO-7010-Symbole verwenden.
% =====================================================================
\BAsetup{
  typ=maschine,
  titel={Dübelfräse Festool DOMINO DF 500 RQ-Plus},
  kurztitel={Dübelfräse},
  taetigkeit={Fräsen von Langlöchern für Dübelverbindungen, Fräserwechsel, Reinigen},
  nummer=BA-DF-01,
  stand=01.10.2026,                   % Datum der Freigabe
  verantwortlich={Name der verantwortlichen Person},
  signalwort=WARNUNG,
  schrift=\footnotesize,
}

\BAkopf

\BAblock{ANWENDUNGSBEREICH}{M002}{%
  \begin{itemize}
    \item Fräsen von Langlöchern für DOMINO-Dübel in Massivholz und Holzwerkstoffe.
    \item Handgeführt, mit beiden Händen; Werkstück fest gespannt.
    \item Benutzung nur nach unterschriebener Einweisung; Betriebsanleitung des Herstellers beachten.
  \end{itemize}}

\BAblock{GEFAHREN FÜR MENSCH UND UMWELT}{W022,W024,W012}{%
  \begin{itemize}
    \item Schnittverletzungen am rotierenden und pendelnden Fräser, der beim Eintauchen vorn aus der Stirnseite austritt und nach dem Ausschalten nachläuft.
    \item Wegrutschen der Fräse, wenn sie beim Eintauchen nicht fest am Werkstück anliegt.
    \item Wegfliegende Späne; Einzug von Haaren, Schmuck, loser Kleidung und Handschuhen; Stromschlag bei beschädigtem Kabel; Lärm.
    \item Gesundheitsschäden durch Staub; Hartholzstaub kann Krebs erzeugen.
  \end{itemize}}

\BAblock{SCHUTZMASSNAHMEN UND VERHALTENSREGELN}{M003,M004,M017,P028}{%
  \begin{itemize}
    \item Vor jeder Benutzung Fräse, Kabel und Fräser prüfen; nur scharfe, unbeschädigte und passende Fräser verwenden.
    \item Gehörschutz und Schutzbrille tragen; Staubmaske bei viel Staub (Hartholz). Anliegende Kleidung, Haare zusammenbinden. \textbf{Beim Fräsen keine Handschuhe} (Einzugsgefahr).
    \item Fräser nur bei \textbf{gezogenem Netzstecker} wechseln.
    \item Werkstück fest spannen, kleine Teile nie mit der Hand halten. Immer Absaugung anschließen (Festool-Sauger, Stellung AUTO).
    \item Fräse mit \textbf{beiden Händen} führen; Hände nie vor die Stirnseite. Erst ansetzen, dann einschalten, dann gleichmäßig eintauchen.
    \item Schalter rastet ein: nach dem Fräsen ausschalten, warten, bis der Fräser steht, erst dann absetzen. Kabel aus dem Arbeitsbereich fernhalten.
  \end{itemize}}

\BAblock{VERHALTEN BEI STÖRUNGEN UND IM GEFAHRENFALL}{W001,M006}{%
  \begin{itemize}
    \item Fräse ausschalten, Stecker ziehen, Betreuer informieren; Schaden und Hergang per Mail an \texttt{fablab-aktive@fablab.fau.de} melden.
    \item Entstehungsbrand: Fräse vom Netz trennen, löschen. \textbf{Feuerwehr: 112}
  \end{itemize}}

\BAblock{VERHALTEN BEI UNFÄLLEN – ERSTE HILFE}{E003}{%
  \begin{itemize}
    \item Fräse ausschalten, Stecker ziehen. Betreuer/Ersthelfer verständigen, Wunden versorgen.
    \item Jeden Unfall melden und im Verbandbuch eintragen.
  \end{itemize}\vspace{1.5mm}
  \textbf{Notruf: 112}\qquad FabLab: 09131\,/\,85\,28013}

\BAletzterblock{INSTANDHALTUNG, ENTSORGUNG}{M021}{%
  \begin{itemize}
    \item Fräserwechsel und Wartung nur durch eingewiesene Personen mit gezogenem Stecker.
    \item Fräse nach jeder Benutzung aussaugen. Elektrische Prüfung nach DGUV Vorschrift 3 gemäß Prüffrist.
  \end{itemize}}
